\uuid{r4kD}
\titre{Primitives usuelles}
\auteur{hanani ; mzari}
\organisation{exo7}
\licence{CC BY-NC-SA 4.0}
\datecreate{2020-01-20}
\niveau{L1}
\module{Analyse}
\chapitre{Calcul d'intégrales}
\sousChapitre{Primitives diverses}
\difficulte{1}
\theme{Primitives, primitives usuelles}
\exo7qcmid{lille-2.primitives.005}

\contenu{
\begin{question}
\qclassification{127.11}
\qtags{motcle=primitives usuelles}
Parmi les affirmations suivantes, cocher celles qui sont vraies :
\begin{answers}
\bad{Une primitive de $\displaystyle f(x)=\mathrm{e}^{-x}\cos x$ sur $\Rr$ est $\displaystyle F(x)=\mathrm{e}^{-x}\cos x$.}
\bad{Une primitive de $\displaystyle f(x)=\mathrm{e}^{-x}\cos x$ sur $\Rr$ est $\displaystyle F(x)=\frac{1}{2}\mathrm{e}^{-x}\left(\cos x+\sin x\right)$.}
\good{Une primitive de $\displaystyle f(x)=\mathrm{e}^{-x}\cos x$ sur $\Rr$ est $\displaystyle F(x)=\frac{1}{2}\mathrm{e}^{-x}\left(\sin x-\cos x\right)$.}
\good{Une primitive de $\displaystyle f(x)=\mathrm{e}^{-x}\cos x$ sur $\Rr$ est $\displaystyle F(x)=\frac{1}{\sqrt{2}}\mathrm{e}^{-x}\sin \left(x-\frac{\pi}{4}\right)$.}
\end{answers}
\begin{explanations}
$\displaystyle \left[\frac{1}{2}\mathrm{e}^{-x}\left(\sin x-\cos x\right)\right]'=\mathrm{e}^{-x}\cos x$ et $\displaystyle \left[\frac{1}{\sqrt{2}}\mathrm{e}^{-x}\sin \left(x-\frac{\pi}{4}\right)\right]'=\mathrm{e}^{-x}\cos x$. Donc $\displaystyle \frac{1}{2}\mathrm{e}^{-x}\left(\sin x-\cos x\right)$ et $\displaystyle \frac{1}{\sqrt{2}}\mathrm{e}^{-x}\sin \left(x-\frac{\pi}{4}\right)$ sont des primitives de $\displaystyle \mathrm{e}^{-x}\cos x$ sur $\Rr$.
\end{explanations}
\end{question}
}
